\item \textbf{ACTIVIDAD} Considere un objeto esférico de radio $r$ y atmósfera a una distancia $d$ del Sol. Supongamos que su emisividad es $e = 1$ para todo tipo de ondas electromagnéticas y que su temperatura es uniforme sobre su superficie. A la distancia de la Tierra $R$ del Sol, la \textit{intensidad de la radiación solar} es $I_s = 1\,370 \text{ W/m}^2$. Esta intensidad varía en $1/d^2$ a distancias distintas de $R$. Un objeto esférico típico absorberá $70.0\%$ de la radiación solar sobre su sección transversal circular $\pi r^2$. (El objeto reflejará aproximadamente $30.0\%$ de la radiación incidente, el objeto parece circular cuando se ve desde el Sol.) Emitirá principalmente radiación infrarroja de toda su superficie $4\pi r^2$. 
(a) Demuestre que la temperatura superficial de equilibrio de un objeto a una distancia $d$ del Sol es
$$ T = \left[ \frac{(0.700)I_s}{4\sigma} \left(\frac{R}{d}\right)^2 \right]^{1/4} \approx (255 \text{ K}) \sqrt{\frac{R}{d}} $$
(b) Use la ecuación en el inciso (a) para determinar la temperatura superficial teórica para los ocho planetas y el planeta enano Plutón, usando la media distancia $d$. (También incluya el planeta enano Ceres, a una distancia de $d = 4.14 \times 10^{11} \text{ m}$ del Sol, para un total de 10 objetos en nuestro sistema solar.) 
(c) Haga una gráfica de barras de las temperaturas encontradas en el inciso (b), así como de su gráfica de barras las temperaturas superficiales teóricas y estimadas, en kelvins, según lo provisto por el Instituto Lunar y Planetario, que se muestran en la tabla adjunta. 
(d) Observe primero a nuestro propio planeta, la Tierra. ¿Existe algún significado, en términos de vida en este planeta, para el hecho de que la temperatura teórica esté por debajo del punto de congelación del agua, mientras que la temperatura medida está por encima de ella? 
(e) La temperatura real de la Tierra se eleva por la absorción atmosférica de la radiación infrarroja emitida desde la superficie. Este efecto a veces se llama \textit{efecto invernadero}. Considere los objetos con las atmósferas más delgadas: Mercurio, Ceres y Plutón. ¿Qué nota sobre la comparación de temperaturas teóricas y medidas para estos planetas? 
(f) Considere los gigantes gaseosos: Júpiter, Saturno, Urano y Neptuno. Estos planetas no tienen superficie sólida; los datos de temperatura se proporcionan para un punto en la atmósfera donde la presión es la misma que en el nivel del mar en la Tierra. ¿Qué nota sobre la comparación de temperaturas teóricas y medidas para estos planetas? 
(g) La discrepancia más clara entre las temperaturas teóricas y medidas en su gráfico es para Venus. ¿Por qué la temperatura medida es mucho más alta que la temperatura teórica? 
(h) ¿Qué puede concluir sobre la atmósfera de Marte a partir de su gráfico?
